% ============================================================
% Appendix: BPM cheat sheets (single-column layout, footnotesize)
% Source: slides/BPM cheat sheet 1..4 (definitions, properties, results, analysis)
% ============================================================

\section*{Appendix: BPM Cheat Sheets}
\addcontentsline{toc}{section}{Appendix: BPM Cheat Sheets}

\footnotesize

% ── BPM cheat sheet 1 — Definitions ─────────────────────────
\subsection*{A.1: Definitions}
\addcontentsline{toc}{subsection}{A.1: Definitions}

\begin{tabularx}{\linewidth}{@{}>{\bfseries\raggedright\arraybackslash}p{0.18\linewidth}L@{}}
\toprule
\rowcolor{rowgray} \multicolumn{2}{c}{\textbf{\color{accent}Definitions}} \\
\midrule
Net & $N = (P, T, F)$ where $P$ is the set of \textbf{places}, $T$ is the set of \textbf{transitions}, $F \subseteq (P \times T) \cup (T \times P)$ is the set of \textbf{arcs} (flow relation). A net is a bipartite, oriented graph whose set of \textbf{nodes} is $P \cup T$. \\
Marking & $M : P \to \mathbb{N}$ such that $M(p)$ is the number of tokens in the place $p$. A marking can be seen as a vector $\mathbb{N}^{|P|}$ (one entry for each place). A place $p$ is \textbf{marked} at $M$ if $M(p) > 0$, otherwise it is \textbf{unmarked}. \\
System & $(N, M_0)$ where $N = (P, T, F)$ is a net and $M_0 : P \to \mathbb{N}$ is the \textbf{initial marking}. \\
S-net / S-system & Every transition has \emph{exactly} one input and one output place ($|{}^\bullet t| = 1 = |t^\bullet|$). \\
T-net / T-system & Every place has \emph{exactly} one input and one output transition ($|{}^\bullet p| = 1 = |p^\bullet|$). \\
Free-choice & Any two transitions must have pre-sets that are either equal or disjoint. \\
Weakly connected & There exists an \textbf{undirected} path between any two nodes. \\
Strongly connected & There exists a \textbf{directed} path between any two nodes. \\
Transition sequence & $\sigma \in T^*$ a possibly empty sequence of transitions. \\
Parikh vector & $\vec\sigma : T \to \mathbb{N}$ such that $\vec\sigma(t)$ counts the number of times $t$ appears in $\sigma$. It is seen as a vector $\mathbb{N}^{|T|}$ (one entry for each transition). \\
Incidence matrix & $N \in \mathbb{Z}^{|P| \times |T|}$ has one row for each place and one column for each transition. $N(p,t) = -1$ if there is an arc from $p$ to $t$ but not vice versa; $N(p,t) = 1$ if there is an arc from $t$ to $p$ but not vice versa; $N(p,t) = 0$ otherwise. \\
S-invariant & $\mathbf{I} \in \mathbb{Q}^{|P|}$ is a vector with one entry for each place such that $\mathbf{I} \cdot N = 0$. It is a map $\mathbf{I} : P \to \mathbb{Q}$ that assigns a weight $\mathbf{I}(p)$ to each place $p$. \\
T-invariant & $\mathbf{J} \in \mathbb{Q}^{|T|}$ is a vector with one entry for each transition such that $N \cdot \mathbf{J} = 0$. It is a map $\mathbf{J} : T \to \mathbb{Q}$ that assigns a multiplicity $\mathbf{J}(t)$ to each transition $t$. \\
Enabling & $M \xrightarrow{t}$ says that every place in the pre-set of $t$ carries a token. In formulas: $\forall p \in {}^\bullet t.\; M(p) \geq 1$, written compactly ${}^\bullet t \subseteq M$ (reading the set ${}^\bullet t$ as the marking with one token per input place). $M \to$ says that some transition is enabled at $M$. \\
Firing & $M \xrightarrow{t} M'$ says that $t$ is enabled at $M$ and its firing leads to $M'$. In formulas: $\bullet t \subseteq M\; \wedge\; M' = M - \bullet t + t\bullet$. $M \xrightarrow{\sigma} M'$ says that $\sigma \in T^*$ is enabled at $M$ and its firing leads to $M'$. In formulas, for $\sigma = t_1 \ldots t_n$ a \textbf{finite} sequence of transitions: $\exists M_1, \ldots, M_n.\; M \xrightarrow{t_1} M_1 \xrightarrow{t_2} \ldots \xrightarrow{t_n} M_n$. \\
Reachable markings & $M \to^* M'$ says that $M'$ is \textbf{reachable} from $M$, i.e.\ that $M \xrightarrow{\sigma} M'$ for some $\sigma \in T^*$. $[M\rangle$ is the set of all markings reachable from $M$. \\
Infinite firings & $M \xrightarrow{\sigma}$ says that $\sigma \in T^\infty$ is enabled at $M$. In formulas, for $\sigma = t_1 t_2 \ldots t_n \ldots$ an \textbf{infinite} sequence of transitions: $\exists M_1, M_2, \ldots, M_n, \ldots.\; M \xrightarrow{t_1} M_1 \xrightarrow{t_2} \ldots \xrightarrow{t_n} M_n \xrightarrow{t_{n+1}} \ldots$ \\
Siphon & $R \subseteq P$ such that $\bullet R \subseteq R\bullet$. \\
Trap & $R \subseteq P$ such that $\bullet R \supseteq R\bullet$. \\
\bottomrule
\end{tabularx}

% ── BPM cheat sheet 2 — Properties ──────────────────────────
\subsection*{A.2: Properties}
\addcontentsline{toc}{subsection}{A.2: Properties}

\begin{tabularx}{\linewidth}{@{}>{\bfseries\raggedright\arraybackslash}p{0.20\linewidth}L@{}}
\toprule
\rowcolor{rowgray} \multicolumn{2}{c}{\textbf{\color{accent}Properties}} \\
\midrule
Liveness & A transition $t$ is \textbf{live} if at any point in time of the computation it is still possible to enable $t$ in the future. In formulas: $\forall M \in [M_0\rangle.\; \exists M' \in [M\rangle.\; M' \xrightarrow{t}$. A system $S$ is \textbf{live} if all its transitions are live. In formulas: $\forall t \in T.\; \forall M \in [M_0\rangle.\; \exists M' \in [M\rangle.\; M' \xrightarrow{t}$. \\
Place-liveness & A place $p$ is \textbf{live} if at any point in time of the computation it is still possible to mark $p$ in the future. In formulas: $\forall M \in [M_0\rangle.\; \exists M' \in [M\rangle.\; M'(p) > 0$. A system $S$ is \textbf{place-live} if all its places are live. In formulas: $\forall p \in P.\; \forall M \in [M_0\rangle.\; \exists M' \in [M\rangle.\; M'(p) > 0$. \\
Dead transition & A transition $t$ is \textbf{dead} at $M$ if it will never be enabled in the future. In formulas: $\forall M' \in [M\rangle.\; M' \not\xrightarrow{t}$. \\
Dead place & A place $p$ is \textbf{dead} at $M$ if it will always remain unmarked. In formulas: $\forall M' \in [M\rangle.\; M'(p) = 0$. \\
Non-live node & A place / transition is \textbf{non-live} if it can become dead in the future, ie.\ if there exists a reachable marking where it is dead. \\
Deadlock freedom & A system is \textbf{deadlock free} if every reachable marking enables some transition. In formulas: $\forall M \in [M_0\rangle.\; \exists t \in T.\; M \xrightarrow{t}$, i.e.\ more simply, $\forall M \in [M_0\rangle.\; M \to$. \\
Deadlock & A marking $M$ is \textbf{deadlock} if it does not enable any transitions. In formulas: $M \not\to$. \\
$k$-boundedness & A place $p$ is $k$-\textbf{bounded} if it will never contain more than $k$ tokens. In formulas: $\forall M \in [M_0\rangle.\; M(p) \le k$. A system is $k$-\textbf{bounded} if all its places are $k$-bounded. In formulas: $\forall p \in P.\; \forall M \in [M_0\rangle.\; M(p) \le k$. \\
Safeness & 1-bounded places and systems are called \textbf{safe}. \\
Boundedness & A place is \textbf{bounded} if there exists $k \in \mathbb{N}$ such that it is $k$-bounded. \\
Positive vector & A vector is \textbf{positive} if all its entries are strictly greater than 0. \\
Semi-positive vector & A vector is \textbf{semi-positive} if it has no negative entries and at least a positive one. \\
Uniform vector & A vector is \textbf{uniform} if its entries are either $0$ or $1$. \\
Support & $\langle V \rangle$ is the set of elements of the vector $V$ that are assigned positive values. \\
Stable set & A set of markings $\mathbf{M}$ is \textbf{stable} if every marking reachable from some marking in $\mathbf{M}$ is also in $\mathbf{M}$. In formulas: $\forall M \in \mathbf{M}.\; \forall M' \in [M\rangle.\; M' \in \mathbf{M}$. \\
Workflow net & Given a net $(P, T, F)$ is a \textbf{workflow net} if: \begin{enumerate}\item There exists a (unique) initial place $i \in P$ such that $\bullet i = \emptyset$. \item There exists a (unique) final place $o \in P$ such that $o\bullet = \emptyset$. \item Every node of the net belongs to an oriented path from $i$ to $o$.\end{enumerate} \\
$N^*$ & Given a workflow net $N$, the net $N^*$ is obtained by adding a reset transition from the final place of $N$ to the initial place of $N$. \\
% NOTE: the source sheet annotates "(Weak soundness = 1 and 2)", but with
% this numbering that would exclude proper completion, contradicting the
% definition verified on deck 18 p.27 (weak soundness = option to complete
% + proper completion, dead tasks allowed). Corrected to "2 and 3":
% suspected off-by-one on the official sheet, exam-relevant.
Soundness \par (Weak soundness $=$ 2 and 3) & Given a workflow net $N$, $N$ is sound if and only if the following conditions hold: \begin{enumerate}\item No dead task: $\forall t \in T.\; \exists M \in [i\rangle.\; M \xrightarrow{t}$ (every transition can be executed at least once). \item Option to complete: $\forall M \in [i\rangle.\; \exists M' \in [M\rangle.\; M'(o) \geq 1$ (at any point in time of the computation there is still the possibility to mark $o$). \item Proper completion: $\forall M \in [i\rangle.\; M(o) > 0 \Rightarrow M = [o]$ (any reachable marking that marks the final place must have exactly one token in $o$ and no token elsewhere). \end{enumerate} \\
\bottomrule
\end{tabularx}

% ── BPM cheat sheet 3 — General results ─────────────────────
\subsection*{A.3: General Results}
\addcontentsline{toc}{subsection}{A.3: General Results}

\begin{tabularx}{\linewidth}{@{}L@{}}
\toprule
\rowcolor{rowgray} \multicolumn{1}{c}{\textbf{\color{accent}General results}} \\
\midrule
\textbf{Proposition}: liveness implies place-liveness. \emph{Consequence}: if a system is not place-live, it cannot be live. \par \textbf{Lemma}: liveness implies deadlock freedom. \par \textbf{Proposition}: place-liveness implies deadlock freedom. \emph{Consequences}: if a system can deadlock, it is neither live nor place-live. \\
\midrule
\textbf{Enabling proposition.} $M \xrightarrow{\sigma}$ if and only if $M \xrightarrow{\sigma'}$ for any (finite) prefix $\sigma'$ of $\sigma$. \\
\midrule
\textbf{Marking equation lemma.} If $M \xrightarrow{\sigma} M'$ then $M' = M + N \cdot \vec\sigma$. \\
\midrule
\textbf{Monotonicity lemma (1).} If $M \xrightarrow{\sigma} M'$ then $M + L \xrightarrow{\sigma} M' + L$ for any $L : P \to \mathbb{N}$. \\
\midrule
\textbf{Monotonicity lemma (2).} If $M \xrightarrow{\sigma}$ then $M + L \xrightarrow{\sigma}$ for any $L : P \to \mathbb{N}$. \\
\midrule
\textbf{Monotonicity corollary.} If $M \xrightarrow{\sigma} M'$ with $M \subseteq M'$ then $M \xrightarrow{\sigma\sigma\ldots}$. \\
\midrule
\textbf{Boundedness lemma.} If a system is bounded and $M_0 \xrightarrow{\sigma} M$ with $M_0 \subseteq M$ then $M = M_0$. \\
\midrule
\textbf{Unboundedness lemma.} A system is not bounded iff $\exists M \in [M_0\rangle.\; \exists M' \in [M\rangle.\; M \subset M'$. \\
\midrule
\textbf{Repetition lemma.} If $M \xrightarrow{\sigma} M'$ and $M \xrightarrow{\sigma\sigma\ldots}$, then $M \subseteq M'$. \\
\midrule
\textbf{Strong connectedness theorem (1).} If a weakly connected system is live and bounded, then it is strongly connected. \emph{Consequence}: if a system is not strongly connected, then it cannot be live and bounded at the same time. \\
\midrule
\textbf{Main Theorem.} A workflow net $N$ is \textbf{sound} if and only if $N^*$ is \textbf{live and bounded}. \\
\midrule
\textbf{Compositionality Theorem.} If $N$ and $N'$ are two sound and safe workflow nets, then $N[N'/t]$ is safe and sound. \\
\midrule
\textbf{Place-liveness $=$ liveness in free-choice systems.} If a free-choice system is place-live, then it is live. \\
\midrule
\textbf{Commoner's Theorem.} A free-choice system is live if and only if every proper siphon includes a marked trap. \\
\midrule
\textbf{Rank Theorem.} A free-choice system $(P, T, F, M_0)$ is live and bounded if and only if \begin{enumerate}\item It has at least one place and one transition. \item It is connected. \item $M_0$ marks every proper siphon. \item It has a positive S-invariant. \item It has a positive T-invariant. \item The rank of the incidence matrix $N$ equals the number of clusters minus 1.\end{enumerate} \\
\midrule
\textbf{Fundamental property of siphons.} Unmarked siphons remain unmarked. \\
\midrule
\textbf{Liveness and siphons proposition.} If a system is live, then any proper siphon is marked. \emph{Corollary}: if a system has an unmarked proper siphon, then it is not live. \\
\midrule
\textbf{Fundamental property of traps.} Marked traps remain marked. \\
\bottomrule
\end{tabularx}

\bigskip
\begin{tabularx}{\linewidth}{@{}L@{}}
\toprule
\rowcolor{rowgray} \multicolumn{1}{c}{\textbf{\color{accent}Invariants}} \\
\midrule
\textbf{Proposition.} $\mathbf{I}$ is an S-invariant if and only if $\forall t \in T.\; \sum_{p \in \bullet t} \mathbf{I}(p) = \sum_{q \in t\bullet} \mathbf{I}(q)$. \\
\midrule
\textbf{Vector space.} Any linear combination of S-invariants is also an S-invariant. \\
\midrule
\textbf{Fundamental property of S-invariants.} $\forall M \in [M_0\rangle.\; \mathbf{I} \cdot M = \mathbf{I} \cdot M_0$. \emph{Corollary}: if $\mathbf{I} \cdot M \neq \mathbf{I} \cdot M_0$, then $M$ is not reachable from $M_0$. \\
\midrule
\textbf{Boundedness theorem for S-invariants} (sufficient condition for boundedness). If a net has a positive S-invariant, then it is bounded (for any initial marking). \emph{Corollary}: any place in the support of a semi-positive S-invariant is bounded (for any initial marking). \\
\midrule
\textbf{Liveness theorem for S-invariants} (necessary condition for liveness). If a system is live, then for any semi-positive S-invariant $\mathbf{I}$ we have $\mathbf{I} \cdot M_0 > 0$. \emph{Corollary}: if we find a semi-positive S-invariant $\mathbf{I}$ such that $\mathbf{I} \cdot M_0 = 0$ the system is not live. \\
\midrule
\textbf{Proposition.} $\mathbf{J}$ is a T-invariant if and only if $\forall p \in P.\; \sum_{t \in \bullet p} \mathbf{J}(t) = \sum_{u \in p\bullet} \mathbf{J}(u)$. \\
\midrule
\textbf{Fundamental property of T-invariants.} Let $M \xrightarrow{\sigma} M'$. Then $\vec\sigma$ is a T-invariant if and only if $M' = M$. \emph{Corollary}: any firing sequence whose Parikh vector is a T-invariant leads back to the starting marking. \\
\midrule
\textbf{Reproduction lemma.} If a system is bounded and $M_0 \xrightarrow{\sigma}$ with $\sigma \in T^\infty$, then there must exists a semi-positive T-invariant $\mathbf{J}$ such that $\langle \mathbf{J} \rangle \subseteq \{t \in \sigma\}$. \\
\midrule
\textbf{Boundedness $+$ liveness theorem for T-invariants.} If a bounded system is live, then it has a positive T-invariant. \emph{Corollary}: every live and bounded system has some $M \in [M_0\rangle$ and $M \xrightarrow{\sigma} M'$ s.t.\ $\{t \in \sigma\} = T$. \\
\midrule
\textbf{Strong connectedness theorem (2).} If a weakly connected system has positive S- and T-invariants, then it is strongly connected. \emph{Consequence}: if a system is not strongly connected, then we cannot find positive S- and T-invariants. \\
\bottomrule
\end{tabularx}

\bigskip
\begin{tabularx}{\linewidth}{@{}L@{}}
\toprule
\rowcolor{rowgray} \multicolumn{1}{c}{\textbf{\color{accent}Diagnostic}} \\
\midrule
\textbf{S-coverability.} If a free-choice system is live and bounded, then it is S-coverable. \emph{Corollary}: free-choice $+$ not S-coverable $\Rightarrow$ not live or not bounded. \\
\midrule
\textbf{S-coverability and workflow-nets.} If $N$ is sound and free-choice, then $N^*$ is S-coverable. \emph{Corollary}: $N$ free-choice $+$ $N^*$ not S-coverable $\Rightarrow N$ is not sound. \\
\midrule
\textbf{Compositionality of sound free-choice nets.} \emph{Lemma}: If a free-choice workflow net is sound, then it is safe. \emph{Proposition}: If $N$ and $N'$ are two sound free-choice workflow nets, then $N[N'/t]$ is sound and free-choice. \\
\midrule
\textbf{Well-structuredness.} If $N$ is sound and well-structured, then $N^*$ is S-coverable. \emph{Corollary}: $N$ well structured $+$ $N^*$ not S-coverable $\Rightarrow N$ is not sound. \\
\midrule
\textbf{Fundamental property of non-live error sequences.} Let $N$ be such that $N^*$ is bounded and has no dead task. Then, $N^*$ is live if and only if $N$ has no non-live sequences. \\
\midrule
\textbf{Fundamental property of unbounded error sequences.} $N^*$ is bounded if and only if $N$ has no unbounded sequences. \\
\bottomrule
\end{tabularx}

% ── BPM cheat sheet 4 — Analysis ────────────────────────────
\subsection*{A.4: Analysis Summary Table}
\addcontentsline{toc}{subsection}{A.4: Analysis Summary Table}

\renewcommand{\arraystretch}{1.35}
\begin{tabularx}{\linewidth}{@{}>{\centering\arraybackslash}p{0.13\linewidth}>{\centering\arraybackslash}p{0.13\linewidth}>{\centering\arraybackslash}p{0.16\linewidth}>{\centering\arraybackslash}p{0.13\linewidth}>{\centering\arraybackslash}p{0.18\linewidth}>{\centering\arraybackslash}p{0.13\linewidth}@{}}
\toprule
\rowcolor{accent!15}
\textbf{S-systems} & \textbf{Strongly connected S-systems} & \textbf{T-systems} & \textbf{Strongly connected T-systems} & \textbf{System} & \textbf{Property} \\
\midrule
$\checkmark$ & $\checkmark$ & $\forall p \in P.\; \exists \gamma \in \Gamma.\; p \in \gamma$ (every place belongs to a circuit) & $\checkmark$ & Positive S-invariant & $\Rightarrow$ boundedness \\
\midrule
$\times$ & $\times$ & $\exists p \in P.\; \forall \gamma \in \Gamma.\; p \notin \gamma$ (there is a place that is not covered by any circuit) & $\times$ & There exist two reachable markings $M$ and $M'$ such that $M \subset M'\; \wedge\; M' \in [M\rangle$ & $\Rightarrow$ unboundedness \\
\midrule
Strongly connected and $M_0(P) > 0$ & $M_0(P) > 0$ & $\forall \gamma \in \Gamma.\; M_0(\gamma) > 0$ & $\forall \gamma \in \Gamma.\; M_0(\gamma) > 0$ & & $\Rightarrow$ liveness \\
\midrule
Not strongly connected or $M_0(P) = 0$ & $M_0(P) = 0$ & $\exists \gamma \in \Gamma.\; M_0(\gamma) = 0$ & $\exists \gamma \in \Gamma.\; M_0(\gamma) = 0$ & There is a semi-positive S-invariant $\mathbf{I}$ such that $\mathbf{I} \cdot M_0 = 0$ & $\Rightarrow$ non-liveness \\
\midrule
 & $M(P) = M_0(P)$ &  &  &  & $\Rightarrow M \in [M_0\rangle$ (Reachability) \\
\midrule
$M(P) \neq M_0(P)$ & $M(P) \neq M_0(P)$ & $\exists \gamma \in \Gamma.\; M(\gamma) \neq M_0(\gamma)$ & $\exists \gamma \in \Gamma.\; M(\gamma) \neq M_0(\gamma)$ & $\exists \mathbf{I}.\; \mathbf{I} \cdot M \neq \mathbf{I} \cdot M_0$ & $\Rightarrow M \notin [M_0\rangle$ (Unreachability) \\
\bottomrule
\end{tabularx}
\renewcommand{\arraystretch}{1.2}

\bigskip
Let $N$ be a workflow net and $N^*$ its corresponding variant with the \emph{reset} transition.
\begin{itemize}
  \item \textbf{Theorem}: $N^*$ is strongly connected.
  \item \textbf{Theorem}: If $N$ is an S-net, then it is sound.
  \item \textbf{Theorem}: If $N^*$ is a T-net whose circuits are all marked, then $N$ is sound.
  \item \textbf{Theorem}: If $N^*$ is a T-net that has at least one unmarked circuit, then $N$ is not sound.
  \item \textbf{Theorem}: If $N$ is free-choice and $N^*$ meets all conditions in the Rank Theorem, then $N$ is sound.
  \item \textbf{Theorem}: If $N$ is free-choice and $N^*$ does not have a positive S-invariant, then $N$ is not sound.
  \item \textbf{Theorem}: If $N$ is free-choice and $N^*$ does not have a positive T-invariant, then $N$ is not sound.
  \item \textbf{Theorem}: If $N$ is free-choice and the rank of $N^*$ is different from the number of its clusters minus 1, then $N$ is not sound.
  \item \textbf{Theorem}: If $N$ is free-choice and the $N^*$ has a proper siphon that does not include the initial place $i$, then $N$ is not sound.
  \item \textbf{Theorem}: If $N$ is free-choice and $N^*$ is not S-coverable, then $N$ is not sound.
  \item \textbf{Theorem}: If $N$ is well-structured and $N^*$ is not S-coverable, then $N$ is not sound.
\end{itemize}

\normalsize
